\begin{enumerate}

\item El análisis BPMN determinó que la priorización debe aplicarse a partir de la etapa ``Interesado'', después del contacto y perfilamiento del lead.

\item El tratamiento diferenciado de la ausencia permitió conservar 49.1\% más leads y 10.0\% más conversiones que el criterio de casos completos.

\item La regresión logística concentró 4.15 veces más conversiones que la selección aleatoria al priorizar el 10.0\% de la cartera.

\item Las 1.6 millones de simulaciones de Monte Carlo demostraron que el modelo superó la selección aleatoria en todos los niveles de capacidad evaluados.

\item El modelo proporcionó un criterio cuantitativo y reproducible para determinar qué leads atender primero según su probabilidad estimada de conversión.

\end{enumerate}